% =====================================================================
\section{Noise}
No quantity in the brain stands still. Neuronal activity fluctuates of its own accord, without external cause, and does so on all time scales at once: from fractions of a second to minutes and longer. The influx of oxidants fluctuates with it. For a system with two states this matters: a random push can carry it across the threshold before the mean load reaches the threshold.
\facts
\begin{itemize}
\item The spectrum of spontaneous brain activity has no preferred scale: the power decays with frequency as a power law over a range from about 0.003 to 500\,Hz.\tS
\item The level of arousal changes in jumps within seconds and returns to its previous value over tens of seconds.\tS
\item Episodes of ``voices'' in patients arise spontaneously, without external cause.\tS
\end{itemize}
\begin{ineq}
This is the problem of escape from a potential well under noise. The mean waiting time for a transition depends exponentially on the ratio of the barrier height to the noise intensity. Near the edge of the window the barrier decays as the distance to the edge to the power $3/2$.\tD\ The exponential dependence means that a modest spread of parameters between cells or between people produces an enormous spread of times --- a heavy-tailed distribution, as for the duration of the prodrome. A calculation of the project also revealed a limitation: fast noise of plausible strength is too weak to carry the system over far from the edge; what matters is slow fluctuation, on scales of minutes and longer.\tD
\end{ineq}
